\subsection{Quantizzazione di un sistema}
Il trucco \`e definire le grandezze attraverso le leggi di conservazione.

Dal formalismo lagrangiano, suppongo che il sistema ha lagrangiana, allora
$$
\delta \mathcal{L} = 0 \qquad \Rightarrow \qquad Q = \frac{\partial \mathcal{L}}{\partial \dot{q}} \, \delta q \; \text{(carica di Noether).}
$$

Infatti l'invarianza della lagrangiana \`e
\begin{align*}
0 = \delta \mathcal{L} &= \frac{\partial \mathcal{L}}{\partial q} \delta q + \frac{\partial \mathcal{L}}{\partial \dot{q}} \delta \dot{q} = \\
&= \frac{d}{dt} \frac{\partial \mathcal{L}}{\partial \dot{q}} \delta q + \frac{\partial \mathcal{L}}{\partial \dot{q}} \delta \dot{q} = \\
&= \frac{d}{dt} \frac{\partial \mathcal{L}}{\partial \dot{q}} \delta q + \frac{\partial \mathcal{L}}{\partial \dot{q}} \frac{d}{dt} \delta q = \\
&= \frac{d}{dt} \left( \frac{\partial \mathcal{L}}{\partial \dot{q}} \delta q \right)
\end{align*}
quindi $ \frac{\partial \mathcal{L}}{\partial \dot{q}} \delta q $ \`e costante. Questo \`e il \textbf{teorema di Noether}.

Per le traslazioni, dal punto di vista quantistico ho un operatore unitario che agisce sui vettori di base
$$
T \ket{q} = \ket{q - \delta} \qquad (T^\dag T = \mathbbm{1})
$$
producendo una nuova base.

Per lo stato $ \ket{\psi} $
\begin{align*}
T \ket{\psi} &= \int_{-\infty}^{+\infty} T \ket{q} \braket{q|\psi} \, dq = \int_{-\infty}^{+\infty} \ket{q - \delta} \braket{q|\psi} \, dq = \\
&= \int_{-\infty}^{+\infty} \ket{q - \delta} \psi(q) \, dq = \int_{-\infty}^{+\infty} \ket{q^\prime} \psi(q^\prime + \delta) \, dq^\prime
\end{align*}
Allora
$$
\bra{q} T \ket{\psi} = \int_{-\infty}^{+\infty} \braket{q|q^\prime} \psi(q^\prime + \delta) \, dq^\prime = \int_{-\infty}^{+\infty} \delta(q^\prime - q) \, \psi(q^\prime + \delta) \, dq^\prime = \psi(q + \delta)
$$
che sviluppata \`e
\begin{align*}
\psi(q + \delta) &= \sum_{k = 0}^\infty \delta^k \frac{1}{k!} \psi^{(k)}(q) = \\
&= \left[ \sum_{k = 0}^\infty \frac{\delta^k}{k!} \left( \frac{d}{dq} \right)^k \right] \psi(q) = \\
&= e^{\delta \frac{d}{dq}} \psi(q)
\end{align*}
Inoltre so che $T_\delta = e^{i \hat{K} \delta}$, dove $\hat{K}^\dag = \hat{K}$, quindi
$$
\psi(q + \delta) = e^{i \delta \left( -i \frac{d}{dq} \right)} \psi(q).
$$

Per trasformazioni infinitesime $ T_\varepsilon = e^{i \varepsilon \hat{K}} = 1 + i \varepsilon \hat{K} + \sigma(\varepsilon^2) $, dove $ \hat{K} $ \`e il \textbf{generatore delle traslazioni}. Allora
$$
\bra{q} 1 + i \delta \hat{K} \ket{\psi} =  \bra{q} e^{i \delta \hat{K}} \ket{\psi} = e^{i \delta \left( -i \frac{d}{dq} \right)} \psi(q) = \psi(q) + i \delta \left( -i \psi^\prime(q) \right)
$$
$$
\boxed{
\braket{q|\hat{K}|\psi} = -i \frac{d}{dq} \psi(q).
}
$$

Ma cosa si conserva di un sistema quando c'\`e invarianza per traslazione?

Supponendo che esista $ S $ tale che $ S^\dag \, S = \mathbbm{1} $ e $ S(t, t_0) |\ket{\psi} = \ket{\psi(t)} $ e quindi $ \braket{\psi(t_0)|\psi(t_0)} = \braket{\psi(t)|\psi(t)} $; avendo invarianza per traslazione per due stati $ \ket{\varphi} $ e $ \ket{\psi} $, segue
$$
\braket{\varphi(t)|\psi(t)} = \braket{\varphi(t)|S(t, t_0)|\psi(t_0)} = \braket{\varphi(t)|T^\dag S(t, t_0) T|\psi(t_0)}
$$
che nel caso infinitesimo \`e
$$
\braket{\varphi(t)|\left( \mathbbm{1} - i \varepsilon \hat{K} \right) \, S(t, t_0) \, \left( \mathbbm{1} + i \varepsilon \hat{K} \right) | \psi(t_0)} =
$$
$$
= \braket{\varphi(t)|S(t, t_0) - i \varepsilon \hat{K} S(t, t_0) + i \varepsilon \hat{K} S(t, t_0)|\psi(t_0)} =
$$
$$
= \braket{\varphi(t)|S(t, t_0)|\psi(t_0)} - i \varepsilon \braket{\varphi(t)|\left[ \hat{K}, S(t, t_0) \right]|\psi(t_0)}
$$
Quindi deve essere $ \left[ \hat{K}, S(t, t_0) \right] = 0 $.

Gli autostati di  $ \hat{K} $ ci sono e $ \hat{K} \ket{k} = k \ket{k} $, quindi
\begin{align*}
0 &= \braket{\varphi(t)|\left[ \hat{K}, S(t, t_0) \right]|\psi(t_0)} = \int_{-\infty}^{+\infty} \braket{\varphi|k} \braket{k|[,]|k^\prime} \braket{k^\prime|\psi} dk \, dk^\prime \\
&= \int_{-\infty}^{+\infty} \braket{\varphi|k} \left( k - k^\prime \right) \braket{k|S(t, t_0)|k^\prime} \braket{k^\prime|\psi} dk \, dk^\prime
\end{align*}
e questo per ogni $\ket{\psi}, \ket{\phi}$, quindi
$$
\text{se } k \neq k^\prime, \text{allora } \braket{k|S(t, t_0)|k^\prime} = 0
$$
cio\`e l'autovalore di $ \hat{K} $ si conserva.

Definisco l'\textbf{operatore impulso} o \textbf{momento}
$$
\hat{p} = \hbar \hat{K}.
$$

\subsubsection{Esempio: commutatore posizione-momento}
$$
\hat{q} \xrightarrow{\text{evoluzione}} T^{-1} \hat{q} \, T = \hat{q} - \delta.
$$
\begin{align*}
T^{-1} \hat{q} \, T \ket{q} &= T^{-1} \hat{q} \ket{q - \delta} = T^{-1} \left( q - \delta \right) \ket{q - \delta} = \\
&= \left( q - \delta \right) T^{-1} \ket{q - \delta} = \left( q - \delta \right) \ket{q}.
\end{align*}
Per trasfromazioni infinitesimali $ T_\varepsilon = \mathbbm{1} + i \varepsilon \hat{K} = \mathbbm{1} + i \varepsilon \frac{1}{\hbar} \hat{p} $, quindi
$$
\left( \mathbbm{1} - i \varepsilon \frac{1}{\hbar} \hat{p} \right) \, \hat{q} \, \left( \mathbbm{1} + i \varepsilon \frac{1}{\hbar} \hat{p} \right) = \hat{q} - \varepsilon
$$
$$
\hat{q} - i \varepsilon \frac{1}{\hbar} \left( \hat{p} \hat{q} - \hat{q} \hat{p} \right) = \hat{q} - \varepsilon
$$
quindi  $ \left[ \hat{p}, \hat{q} \right] = -i \hbar $.

\subsubsection{Esempio: ancora sull'operatore traslazione}
$$
T = \begin{cases}
\mathbbm{1} + i \frac{\varepsilon}{\hbar} \hat{p} = T_\varepsilon \text{ infinitesimo} \\
e^{i \frac{\varepsilon}{\hbar} \hat{p}} = T_\varepsilon
\end{cases}
$$
\begin{align*}
\braket{k'|T_\varepsilon|k} &= \braket{k'|\mathbbm{1} + i \frac{\varepsilon}{\hbar} \hat{p}|k} = \delta(k' - k) (1 + i \varepsilon k), \\
\braket{k'|T_\varepsilon|k} &= e^{i \varepsilon k} \delta(k' - k), \\
\braket{k'|\hat{q}|k} &= - i \frac{d}{dk} \delta(k - k') = i \frac{d}{dk'} \delta(k' - k).
\end{align*}